\DocumentMetadata{
  pdfversion=2.0,
  pdfstandard=ua-2,
  lang=en-US,
  tagging=on,
  tagging-setup={math/setup=mathml-SE,tabsorder=structure}
}
\documentclass[11pt,letterpaper]{article}
\usepackage[margin=0.68in,includefoot]{geometry}
\usepackage{newcomputermodern}
\usepackage{amsmath,amscd,mathtools}
\providecommand{\square}{\mdlgwhtsquare}
\usepackage[hidelinks]{hyperref}
\hypersetup{
  pdftitle={Numerical Analysis PhD Exam, September 2003},
  pdfauthor={University of Florida Department of Mathematics},
  pdfsubject={Graduate mathematics examination},
  pdfdisplaydoctitle=true,
  bookmarksopen=true
}
\setcounter{secnumdepth}{0}
\setlength{\parindent}{0pt}
\setlength{\parskip}{0.28em}
\setlength{\emergencystretch}{3em}
\setlength{\leftmargini}{2.15em}
\setlength{\leftmarginii}{2.65em}
\setlength{\leftmarginiii}{3em}
\renewcommand{\labelenumi}{\textbf{\theenumi.}}
\pagestyle{empty}
\raggedbottom
\allowdisplaybreaks
\newenvironment{examproblems}[1]{%
  \begin{enumerate}%
  \setcounter{enumi}{#1}%
  \setlength{\topsep}{0.35em}%
  \setlength{\partopsep}{0pt}%
  \setlength{\itemsep}{0.48em}%
  \setlength{\parsep}{0.18em}%
}{\end{enumerate}}
\newenvironment{examsubparts}{%
  \begin{enumerate}%
  \setlength{\topsep}{0.18em}%
  \setlength{\partopsep}{0pt}%
  \setlength{\itemsep}{0.16em}%
  \setlength{\parsep}{0.1em}%
}{\end{enumerate}}
\newenvironment{examinstructions}{%
  \par\begingroup\itshape
}{%
  \par\endgroup\vspace{0.18em}
}
\makeatletter
\renewcommand\section{\@startsection{section}{1}{0pt}%
  {-1.25ex plus -.25ex minus -.1ex}%
  {0.55ex plus .1ex}%
  {\normalfont\large\bfseries}}
\renewcommand{\@maketitle}{%
  \newpage\null\vskip -1.2em%
  {\footnotesize\bfseries
    \href{https://www.ufl.edu/}{University of Florida}, \href{https://math.ufl.edu/}{Department of Mathematics}\par}%
  \vskip 0.18em%
  \tagstructbegin{tag=Title}%
    {\LARGE\bfseries\href{https://gma.math.ufl.edu/exams/numerical-analysis-phd/}{Numerical Analysis PhD Exam}, September 2003\par}%
  \tagstructend%
  \vskip 0.35em%
  \tagmcbegin{artifact}%
    \hrule height 1pt%
  \tagmcend%
  \vskip 0.55em%
}
\makeatother
\title{Numerical Analysis PhD Exam, September 2003}
\author{}
\date{}
\begin{document}
\maketitle
\thispagestyle{empty}
\begin{examinstructions}
Do any 8 of the following 10 problems:
\end{examinstructions}
\begin{examproblems}{0}
\item
This problem has three parts.
\begin{examsubparts}
\item[\textnormal{(a)}]
Let~\(V\) be the space of all complex \(m\times n\) matrices endowed with any norm. Prove that its subset of full rank matrices is dense in~\(V\).
\item[\textnormal{(b)}]
Suppose~\(A\) and~\(B\) are square matrices.
\begin{examsubparts}
\item[\textnormal{i.}]
If~\(A\) is invertible, show that \(AB\) and \(BA\) are similar matrices (so their spectra coincide).
\item[\textnormal{ii.}]
How will you use Problem (1a) to conclude that the spectra of \(AB\) and \(BA\) coincide when nothing is known about the invertibility of~\(A\) or~\(B\)?
\end{examsubparts}
\item[\textnormal{(c)}]
For all rectangular matrices \(A\in\mathbb C^{m\times n}\), \(B\in\mathbb C^{n\times m}\) (\(m\ne n\)), does the spectrum of \(AB\) and \(BA\) coincide?
\end{examsubparts}
\item
Let \(A\in\mathbb C^{m\times m}\), and let \(E\in\mathbb C^{m\times m}\) be of rank \(k<m\). Suppose the matrix \(M=A+E\) is invertible. Consider GMRES applied to the preconditioned system 
\[
M^{-1}Ax=M^{-1}b.
\]
 Give a tight bound, independent of the starting guess, on the number of GMRES iterations needed to compute the exact solution~\(x\) (ignoring roundoff). Your answer should be less than~\(m\). Give clear reasons.
\item
Let~\(U\) be the upper triangular matrix obtained by applying Gaussian elimination with partial pivoting (the “\(PA=LU\)” factorization) to a matrix \(A\in\mathbb C^{m\times m}\). Show that the growth factor 
\[
\rho=\frac{\max_{i,j}|U_{ij}|}{\max_{i,j}|A_{ij}|}
\]
 satisfies 
\[
\rho\le 2^{m-1}.
\]
 Exhibit a matrix~\(A\) for which the growth factor is \(2^{m-1}\).
\item
This problem has two parts.
\begin{examsubparts}
\item[\textnormal{(a)}]
Suppose \(A=QR\) is a QR factorization of a symmetric tridiagonal~\(A\). Which entries of~\(R\) and~\(Q\) are zero, in general?
\item[\textnormal{(b)}]
State the (unshifted) QR iteration for eigenvalues. Show that if one iterate of the QR iteration is symmetric and tridiagonal, then the next one is also symmetric and tridiagonal.
\end{examsubparts}
\item
Given \(\mathbf v\) and \(\mathbf w\in\mathbb R^n\), show that 
\[
\mathbf v\mathbf w^{\mathsf T}+\mathbf w\mathbf v^{\mathsf T}
=\frac{\lVert\mathbf v\rVert\lVert\mathbf w\rVert}{2}
\left(\mathbf a\mathbf a^{\mathsf T}-\mathbf b\mathbf b^{\mathsf T}\right)
\]
 
\[
=\begin{bmatrix}
\dfrac{\mathbf a}{\lVert\mathbf a\rVert} & \big| & \dfrac{\mathbf b}{\lVert\mathbf b\rVert}
\end{bmatrix}
\begin{bmatrix}
\mathbf v^{\mathsf T}\mathbf w+\lVert\mathbf v\rVert\lVert\mathbf w\rVert & 0\\
0 & \mathbf v^{\mathsf T}\mathbf w-\lVert\mathbf v\rVert\lVert\mathbf w\rVert
\end{bmatrix}
\begin{bmatrix}
\dfrac{\mathbf a}{\lVert\mathbf a\rVert} & \big| & \dfrac{\mathbf b}{\lVert\mathbf b\rVert}
\end{bmatrix}^{\mathsf T},
\]
 where 
\[
\mathbf a=\frac{\mathbf v}{\lVert\mathbf v\rVert}+\frac{\mathbf w}{\lVert\mathbf w\rVert}
\quad\text{and}\quad
\mathbf b=\frac{\mathbf v}{\lVert\mathbf v\rVert}-\frac{\mathbf w}{\lVert\mathbf w\rVert}.
\]
 What is \(\mathbf a^{\mathsf T}\mathbf b\)? What are the eigenvalues of \(\mathbf v\mathbf w^{\mathsf T}+\mathbf w\mathbf v^{\mathsf T}\)?
\item
Show that \(x=1+\tan^{-1}(x)\) has a solution~\(\alpha\). Find an interval \([a,b]\) containing~\(\alpha\) such that for every \(x_0\in[a,b]\) the iteration 
\[
x_{n+1}=1+\tan^{-1}(x_n)
\]
 will converge to~\(\alpha\).
\item
Let \(x_0,\ldots,x_n\) be distinct real points. Let 
\[
\{l_j(x)\}_{j=0,\ldots,n}
\]
 be the Lagrange basis functions for these points. Prove that 
\[
\sum_{j=0}^n l_j(x)=1
\]
 for all~\(x\).
\item
Suppose you want to solve numerically the ode 
\[
y'=f(x,y),\qquad y(0)=y_0.
\]
\begin{examsubparts}
\item[\textnormal{(a)}]
By applying the midpoint rule to the integral in the equation 
\[
y(x_{n+1})=y(x_{n-1})+\int_{x_{n-1}}^{x_{n+1}}f(s,y(s))\,ds,
\]
 derive the midpoint method for odes.
\item[\textnormal{(b)}]
Find the order of the local truncation error for this method.
\item[\textnormal{(c)}]
Determine if this method is stable or not.
\item[\textnormal{(d)}]
Would this method be a good choice for a stiff system such as \(y'=-100y\)? Why or why not?
\end{examsubparts}
\item
This problem has two parts.
\begin{examsubparts}
\item[\textnormal{(a)}]
Find the linear least squares approximation to \(f(x)=x^2\) on the interval \([0,2]\) by optimizing the least squares error over~\(a\) and~\(b\) in \(r^*(x)=ax+b\).
\item[\textnormal{(b)}]
Find the first two orthonormal polynomials \(\phi_0(x),\phi_1(x)\) on \([0,2]\) (weight function \(w(x)=1\)). Use these to find the linear least squares approx to \(f(x)=x^2\), and check this result agrees with the previous problem.
\end{examsubparts}
\item
Derive the one point Gaussian quadrature formula for 
\[
I=\int_0^1 xf(x)\,dx\approx\sum_{j=1}^n w_jf(x_j)
\]
 with weight function \(w(x)=x\). Find the system to solve for the two point Gaussian quadrature formula. (You do not need to solve the system)
\end{examproblems}
\end{document}
